\section{Related Work}\label{sec:related}
\subsection{Few-Shot Industrial Anomaly Detection}
Classical industrial AD typically assumed access to normal-only training data and identified deviations at test time. Patch distribution, self-supervised, and reconstruction-discriminative approaches such as PaDiM, CutPaste, and DRAEM provided important context for this literature~\cite{padim2021,cutpaste2021,draem2021}. PatchCore~\cite{roth2022towards} established memory-bank patch retrieval as a robust baseline by storing representative normal features and scoring test patches using nearest-neighbor distance. AnomalyDINO~\cite{damm2025anomalydino} extended this paradigm to few-shot inspection with frozen DINOv2 representations. Subsequent work introduced registration-based adaptation (RegAD~\cite{regad2022}), graph-structured memory (GraphCore~\cite{graphcore2023}), fast feature reconstruction (FastRecon~\cite{fastrecon2023}), vision-language zero- and few-shot scoring (WinCLIP~\cite{winclip2023}, CLIP-SAM collaboration~\cite{clipsam2025}, and cross-modal prompt regularization~\cite{uniadprompt2026}), and in-context residual learning (InCTRL~\cite{inctrl2024}). These methods were effective rankers, but their reference state could remain nontrivial at inference time and the raw score scale did not represent a calibrated probability. Calibration under distribution shift was also known to degrade for deep models generally~\cite{ovadia2019trust}. Our work is complementary: it does not compete on ranking but instead examines how few-shot detectors should expose decision reliability.

\subsection{Frozen-Feature Subspace Detectors}
Subspace approaches replaced explicit patch retrieval with a compact normal subspace. SubspaceAD~\cite{subspacead2026} demonstrated that frozen DINOv2 patch features and PCA or subspace residuals were already effective for few-shot anomaly ranking, establishing the ranking mechanism as prior art. Our contribution begins after the residual is computed and examines the reliability claims that a low-storage subspace detector can support under few-shot and distribution-shift conditions.

\subsection{Calibration and Decision Reliability in Few-Shot AD}
Calibration measured agreement between predicted confidence and empirical correctness~\cite{platt1999,guo2017calibration}. Standard post-hoc methods included Platt scaling and temperature scaling~\cite{platt1999,guo2017calibration}. Histogram binning and isotonic regression provided nonparametric alternatives~\cite{zadrozny2001obtaining,zadrozny2002transforming}. Sparse labels, open-ended anomaly types, and category-dependent scores made this objective especially difficult in anomaly detection. Prior reliability analyses of few-shot AD based on DINOv2 also documented calibration and perturbation sensitivity~\cite{khan2025calibration}. These findings motivated our separation of ranking quality from an operational question: whether a declared alarm rule maintains its requested false-alarm behavior under a specified shift. A calibrated probability summary alone would not constitute such a risk guarantee.

\subsection{Conformal Prediction and Risk-Controlled Calibration}
Conformal prediction transformed nonconformity scores into p-values under exchangeability assumptions~\cite{vovk2005,shafer2008tutorial,angelopoulos2021gentle}. This approach was particularly suitable for normal-only few-shot AD because normal support images could define a reference distribution without requiring real anomalies. Conformal anomaly detection was already well established: leave-one-out, bootstrap, and cross-conformal anomaly detectors had been directly analyzed~\cite{hennhofer2024loio}; weighted conformal detection in low-data regimes had examined resolution and effective-sample-size effects~\cite{hennhofer2026weighted}; general nonconformity toolkits had been developed~\cite{hennhofer2026nonconform}; and few-shot conformal prediction with auxiliary tasks had addressed small target calibration sets~\cite{fisch2021fewshot}. Conformal AD, LOIO calibration, auxiliary-task transfer, and weighted conformal p-values were therefore established mechanisms.

Risk-controlling calibration also preceded our work. Learn then Test cast finite-sample calibration as a multiple-testing problem over candidate configurations~\cite{angelopoulos2021learntest}, while conformal risk control selected a parameter to control the expectation of a monotone loss~\cite{angelopoulos2024crc}. In anomaly detection, PAC-Wrap provided semi-supervised probably approximately correct (PAC) bounds for false-positive and false-negative rates around a base detector~\cite{li2022pacwrap}. These works established the general principles of candidate testing, held-out risk assessment, and one-sided error bounds; we do not claim those statistical primitives as new. CRESS instead instantiates them for normal-only cross-category thresholding by assigning disjoint source categories to reference, proposal, and certification roles, distinguishing selected-mixture image risk from marginal new-category risk, and auditing whether the available number of independent categories can support a nonzero threshold.

Structure-aware conformal adaptation had also been studied outside anomaly detection. Lam and Anh proposed HeAD-CP, which adapted graph-score diffusion using label-free local-homophily estimates and showed that uniform diffusion could impair prediction-set efficiency on heterophilic graphs while preserving marginal coverage~\cite{lam2026headcp}. This result supported the broader need to audit reliability transformations under structural shift, but its setting concerned graph neural network (GNN) node-classification sets rather than few-shot industrial anomaly alarms or cross-category source certification.

The deterministic source-view choices were conceptually informed by the shared-versus-specialized selection motif in Shape-Adapting Gated Experts (SAGE)~\cite{Thai_2026_CVPR}. That work concerned learned expert selection rather than certification of anomaly thresholds. CRESS uses no learned expert gate: its source-view modes are predeclared protocol choices, and its reference, proposal, and certification partitions have disjoint statistical roles. The citation therefore records design inspiration without attributing CRESS validity or novelty to the SAGE architecture.

The present contribution begins at the independent-category requirement of a simultaneous source certificate, with the target-only grid $1/(k{+}1)$ serving as a motivating finite-resolution constraint. This work establishes the Hoeffding-specific and distribution-free all-zero category-feasibility bounds, empirically audits corruption-shift deviations, and evaluates source-assisted candidate thresholds without target anomaly labels. CRESS does not introduce a new nonconformity score or generic confidence bound; it specifies how source categories are separated into reference, proposal, and certification roles and audits whether the resulting source certificate can be nonvacuous at the available category count.
